\chapter{Deferred Proofs}

\vfill
\localtableofcontents
\vfill
\newpage

\section[Proof of Proposition 2.3]{Proof of \Cref{th:Kalman}}\label{sec:proof_Kalman}

\textbf{Prediction.} Let $\mathbb{V}[X]$ denote the variance of the random variable $X$:
\begin{equation}
	\mathbb{V}[X]\triangleq\mathbb{E}\left[(X-\mathbb{E}[X])(X-\mathbb{E}[X])^\top\right]
\end{equation}

Since $x_k=F_k\:x_{k-1}+b_k^{(e)}+w_k$ is an affine transformation of two independent Gaussian random variables $x_{k-1}|\mathbf{z}_{1:k-1}\sim\mathcal{N}\left(\hat{x}_{k-1|k-1},P_{k-1|k-1}\right)$ and $w_k\sim\mathcal{N}(0,Q_k)$, the predicted state $x_k|\mathbf{z}_{1:k-1}$ is also Gaussian, with:
\begin{equation}
	\begin{aligned}
	\hat{x}_{k|k-1}&\triangleq\mathbb{E}[x_k|\mathbf{z}_{1:k-1}]\\
	&=F_k\:\mathbb{E}[x_{k-1}|\mathbf{z}_{1:k-1}]+b_k^{(e)}+\mathbb{E}[w_k|\mathbf{z}_{1:k-1}]\\
	&=F_k\:\hat{x}_{k-1|k-1}+b_k^{(e)}
	\end{aligned}
\end{equation}
and:
\begin{equation}
	\begin{aligned}
	P_{k|k-1}&\triangleq\mathbb{V}[x_k|\mathbf{z}_{1:k-1}]\\
	&=F_k\:\mathbb{V}[x_{k-1}|\mathbf{z}_{1:k-1}]\:F_k^\top+\mathbb{V}[w_k|\mathbf{z}_{1:k-1}]\\
	&=F_k\:P_{k-1|k-1}\:F_k^\top+Q_k
	\end{aligned}
\end{equation}

The predicted covariance is thus SPD as the sum of an SPSD matrix and the SPD matrix $Q_k$.

\textbf{Correction.} The joint distribution of $(x_k,z_k)$ given $\mathbf{z}_{1:k-1}$ is also Gaussian, since $z_k=H_k\:x_k+b_k^{(o)}+n_k$ is an affine transformation of the independent Gaussian variables $x_k|\mathbf{z}_{1:k-1}\sim\mathcal{N}\left(\hat{x}_{k|k-1},P_{k|k-1}\right)$ and $n_k\sim\mathcal{N}(0,R_k)$. It comes that:
\begin{equation}
	\left.\begin{bmatrix}
	x_k\\z_k
	\end{bmatrix}\right|
	\mathbf{z}_{1:k-1}\sim\mathcal{N}
	\left(\begin{bmatrix}
	\hat{x}_{k|k-1}\\H_k\hat{x}_{k|k-1}+b_k^{(o)}
	\end{bmatrix},
	\begin{bmatrix}
	P_{k|k-1}&P_{k|k-1}H_k^\top\\
	H_k P_{k|k-1}&S_k
	\end{bmatrix}\right)
\end{equation}
where:
\begin{equation}
	S_k\triangleq H_k\:P_{k|k-1}\:H_k^\top+R_k
\end{equation}
is called the innovation covariance, and is the covariance associated with the innovation vector:
\begin{equation}
	y_k\triangleq z_k-\left(H_k\:\hat{x}_{k|k-1}+b_k^{(o)}\right)
\end{equation}

Since $R_k\succ0$, the independent pair $(x_k,n_k)$ has the covariance:
\begin{equation}
	\begin{bmatrix}
	P_{k|k-1}&0\\0&R_k
	\end{bmatrix}\succ0
\end{equation}

Its image under the invertible linear map $(x,n)\mapsto(x,H_k\:x+n)$ therefore has an SPD covariance as well. The Gaussian conditioning formula (see \Cref{th:Gaussian_conditioning}) can thus be applied to give the posterior distribution $x_k|\mathbf{z}_{1:k}$ is Gaussian, with:
\begin{align}
	\hat{x}_{k|k}&=\hat{x}_{k|k-1}+K_k\:y_k\\
	P_{k|k}&=(I-K_k\:H_k)\:P_{k|k-1}
\end{align}
where:
\begin{equation}
	K_k\triangleq P_{k|k-1}\:H_k^\top\:S_k^{-1}
\end{equation}
is the Kalman gain.

Finally, \Cref{th:Gaussian_conditioning} proves that the Schur complement of the SPD joint covariance:
\begin{equation}
	P_{k|k-1}-\left(P_{k|k-1}\:H_k^\top\right)S_k^{-1}\left(H_k\:P_{k|k-1}\right)=P_{k|k-1}-K_k\:H_k\:P_{k|k-1}=P_{k|k}
\end{equation}
is itself SPD, so $P_{k|k}\succ0$.

\section[Proof of Proposition 2.10]{Proof of \Cref{th:RBPF_decorrelation}}\label{sec:proof_RBPF}

As a linear combination of the centred noises $w_k^l$ and $w_k^n$, $\overline{w}_k^l$ is centred. Its covariance develops as:
\begin{equation}
	\begin{aligned}
	\overline{Q}_k^l&=\mathbb{E}\left[\overline{w}_k^l\:{\overline{w}_k^l}^\top\right]\\
	&=\mathbb{E}\left[(w_k^l-Q_k^{nl}\:(Q_k^n)^{-1}\:w_k^n)\:(w_k^l-Q_k^{nl}\:(Q_k^n)^{-1}\:w_k^n)^\top\right]\\
	&=\mathbb{E}\left[w_k^l\:(w_k^l)^\top\right]-\mathbb{E}\left[w_k^l\:(Q_k^{nl}\:(Q_k^n)^{-1}\:w_k^n)^\top\right]-\mathbb{E}\left[(Q_k^{nl}\:(Q_k^n)^{-1}\:w_k^n)\:(w_k^l)^\top\right]\\
	&\qquad+\mathbb{E}\left[(Q_k^{nl}\:(Q_k^n)^{-1}\:w_k^n)\:(Q_k^{nl}\:(Q_k^n)^{-1}\:w_k^n)^\top\right]\\
	&=\mathbb{E}\left[w_k^l\:(w_k^l)^\top\right]-\mathbb{E}\left[w_k^l\:(w_k^n)^\top\right](Q_k^n)^{-\top}\:(Q_k^{nl})^\top-Q_k^{nl}\:(Q_k^n)^{-1}\:\mathbb{E}\left[w_k^n\:(w_k^l)^\top\right]\\
	&\qquad+Q_k^{nl}\:(Q_k^n)^{-1}\:\mathbb{E}\left[w_k^n\:(w_k^n)^\top\right]\:(Q_k^n)^{-\top}\:(Q_k^{nl})^\top\\
	&=Q_k^l-Q_k^{nl}\:(Q_k^n)^{-\top}\:(Q_k^{nl})^\top-Q_k^{nl}\:(Q_k^n)^{-1}\:(Q_k^{nl})^\top+Q_k^{nl}\:(Q_k^n)^{-1}\:Q_k^n\:(Q_k^n)^{-\top}\:(Q_k^{nl})^\top\\
	&=Q_k^l-Q_k^{nl}\:(Q_k^n)^{-1}\:(Q_k^{nl})^\top-Q_k^{nl}\:(Q_k^n)^{-1}\:(Q_k^{nl})^\top+Q_k^{nl}\:(Q_k^n)^{-1}\:(Q_k^{nl})^\top\\
	&=Q_k^l-Q_k^{nl}\:(Q_k^n)^{-1}\:(Q_k^{nl})^\top
	\end{aligned}
\end{equation}

Finally, computing:
\begin{equation}
	\begin{aligned}
	\mathbb{E}\left[\overline{w}_k^l\:(w_k^n)^\top\right]&=\mathbb{E}\left[(w_k^l-Q_k^{nl}\:(Q_k^n)^{-1}\:w_k^n)\:(w_k^n)^\top\right]\\
	&=\mathbb{E}\left[w_k^l\:(w_k^n)^\top\right]-Q_k^{nl}\:(Q_k^n)^{-1}\:\mathbb{E}\left[w_k^n\:(w_k^n)^\top\right]\\
	&=Q_k^{nl}-Q_k^{nl}\:(Q_k^n)^{-1}\:Q_k^n\\
	&=0
	\end{aligned}
\end{equation}
shows that $\overline{w}_k^l$ and $w_k^n$ are uncorrelated.

\section[Proof of Proposition 3.2]{Proof of \Cref{th:kernel_gram_equivalence}}\label{sec:proof_definite_kernels}

\textbf{Forward implication.} Let $n\in\mathbb{N}^*$ and $\mathbf{x}\in\mathcal{U}^n$. It comes from the symmetry of $k$ that:
\begin{equation}
	\begin{aligned}
	k(\mathbf{x},\mathbf{x})&=\begin{bmatrix}
	k(x_1,x_1)&\cdots&k(x_1,x_n)\\\vdots&\ddots&\vdots\\k(x_n,x_1)&\cdots&k(x_n,x_n)
	\end{bmatrix}\\
	&=\begin{bmatrix}
	k(x_1,x_1)^\top&\cdots&k(x_n,x_1)^\top\\\vdots&\ddots&\vdots\\k(x_1,x_n)^\top&\cdots&k(x_n,x_n)^\top
	\end{bmatrix}\\
	&=\begin{bmatrix}
	k(x_1,x_1)&\cdots&k(x_1,x_n)\\\vdots&\ddots&\vdots\\k(x_n,x_1)&\cdots&k(x_n,x_n)
	\end{bmatrix}^\top\\
	&=k(\mathbf{x},\mathbf{x})^\top
	\end{aligned}
\end{equation}

Let $a\in\mathbb{R}^{nd_z}$. We decompose $a=[a_1^\top,\dots,a_n^\top]^\top$ with $a_i\in\mathbb{R}^{d_z},\:i\in[\![1,n]\!]$ to obtain:
\begin{equation}
	\begin{aligned}
	a^\top k(\mathbf{x},\mathbf{x})\:a&=\begin{bmatrix}a_1^\top&\cdots&a_n^\top\end{bmatrix}\begin{bmatrix}k(x_1,x_1)&\cdots&k(x_1,x_n)\\\vdots&\ddots&\vdots\\k(x_n,x_1)&\cdots&k(x_n,x_n)\end{bmatrix}\begin{bmatrix}a_1\\\vdots\\a_n\end{bmatrix}\\
	&=\sum_{i=1}^n\sum_{j=1}^na_i^\top k(x_i,x_j)\:a_j\\
	&\geq0
	\end{aligned}
\end{equation}

\textbf{Reverse implication.} Let $x,x'\in\mathcal{U}$. By symmetry of $k((x,x'),(x,x'))$, it comes that:
\begin{equation}
	\begin{aligned}
	\begin{bmatrix}
	k(x,x)&k(x,x')\\k(x',x)&k(x',x')
	\end{bmatrix}&=\begin{bmatrix}
	k(x,x)&k(x,x')\\k(x',x)&k(x',x')
	\end{bmatrix}^\top\\
	&=\begin{bmatrix}
	k(x,x)^\top&k(x',x)^\top\\k(x,x')^\top&k(x',x')^\top
	\end{bmatrix}
	\end{aligned}
\end{equation}

In particular, the upper-right block gives:
\begin{equation}
	k(x,x')=k(x',x)^\top
\end{equation}

Let $n\in\mathbb{N}^*$, $\mathbf{x}\in\mathcal{U}^n$, and $a_1,\dots,a_n\in\mathbb{R}^{d_z}$, and set $a\triangleq[a_1^\top,\dots,a_n^\top]^\top\in\mathbb{R}^{nd_z}$. Then:
\begin{equation}
	\begin{aligned}
	\sum_{i=1}^n\sum_{j=1}^na_i^\top k(x_i,x_j)\:a_j&=\begin{bmatrix}a_1^\top&\cdots&a_n^\top\end{bmatrix}\begin{bmatrix}k(x_1,x_1)&\cdots&k(x_1,x_n)\\\vdots&\ddots&\vdots\\k(x_n,x_1)&\cdots&k(x_n,x_n)\end{bmatrix}\begin{bmatrix}a_1\\\vdots\\a_n\end{bmatrix}\\
	&=a^\top k(\mathbf{x},\mathbf{x})\:a\\
	&\geq0
	\end{aligned}
\end{equation}

\textbf{Strict version.} The same computations hold with strict inequalities, provided the inputs are pairwise distinct and the coefficients not all zero.

\section[Proof of Proposition 3.3]{Proof of \Cref{th:Matérn_properties}}\label{sec:proof_Matérn}

Assume that $\nu\in\mathbb{R}_+^*$ and let $k_\nu$ denote such a function. We also let:
\begin{equation}
	g:\left\{\begin{aligned}
	\mathbb{R}_+&\to\mathbb{R}_+\\
	r&\mapsto\sigma_f^2\:\frac{2^{1-\nu}}{\Gamma(\nu)}r^\nu K_\nu(r)
	\end{aligned}\right.
	\quad\text{and}\quad
	\tilde{r}:\left\{\begin{aligned}
	\mathbb{R}^{d_x}\times\mathbb{R}^{d_x}&\to\mathbb{R}_+\\
	x,x'&\mapsto\frac{\sqrt{2\nu}}{l}\|x-x'\|
	\end{aligned}\right.
\end{equation}

$\tilde{r}$ and $K_\nu$ are continuous on $\mathbb{R}^{d_x}\times\mathbb{R}^{d_x}$ and $\mathbb{R}_+^*$ respectively, and the latter admits the following asymptotic expansion as $r\to0^+$:
\begin{equation}
	K_\nu(r)=2^{\nu-1}\Gamma(\nu)\:r^{-\nu}+o\left(r^{-\nu}\right)
\end{equation}

Therefore, $f$ can be continuously extended to $0$ by setting $f(0)=\sigma_f^2$, and $k_\nu$ is thus continuous on the entirety of $\mathbb{R}^{d_x}\times\mathbb{R}^{d_x}$.

Let $x,x'\in\mathbb{R}^{d_x}$. By symmetry of the norm, $\tilde{r}(x,x')=\tilde{r}(x',x)$, and thus:
\begin{equation}
	k_\nu=g\circ\tilde{r}
\end{equation}
is symmetric, with $\circ$ being the composition operator.

Furthermore, Bochner's theorem \cite{bib:Akhiezer} states that a continuous integrable function with a strictly positive Fourier transform is SPD. It can be shown \cite{bib:Matérn-correlation} that the Fourier transform of $k_\nu$ on $\mathbb{R}^{d_x}$ is:
\begin{equation}
	S_\nu:\left\{\begin{aligned}
	\mathbb{R}^{d_x}&\to\mathbb{R}\\
	\omega&\mapsto C\:\left(\frac{2\nu}{l^2}+\|\omega\|^2\right)^{-(\nu+d_x/2)}
	\end{aligned}\right.
\end{equation}
where $C>0$ is a normalising constant which only depends on $d_x$ and the parameters of $k_\nu$. Since $\nu>0$, $l>0$, it comes that:
\begin{equation}
	\forall\omega\in\mathbb{R}^{d_x},\:S_\nu(\omega)>0
\end{equation}
and thus $k_\nu$ is a SPD kernel. By definition, it is also a radial kernel.

The case $\nu=+\infty$ is more direct. The squared exponential kernel is continuous on the whole of $\mathbb{R}_+$ without any extension, and symmetric by the same composition argument. Its Fourier transform on $\mathbb{R}^{d_x}$ is:
\begin{equation}
	S_\nu:\left\{\begin{aligned}
	\mathbb{R}^{d_x}&\to\mathbb{R}\\
	\omega&\mapsto C\:\exp\left(-\frac{l^2\|\omega\|^2}{2}\right)
	\end{aligned}\right.
\end{equation}
which is again strictly positive, so Bochner's theorem applies as above.

\section[Proof of Proposition 3.8]{Proof of \Cref{th:ordinary_kriging}}\label{sec:proof_OK}

Let:
\begin{equation}
	\mathcal{L}(W,\xi)\triangleq\operatorname{tr}\left(\mathbb{V}\left[W^\top\tilde{\mathbf{z}}-z\right]\right)+2\:\operatorname{tr}\left(\xi^\top\left(C^\top W-I_{d_z}\right)\right)
\end{equation}
be the Lagrangian associated with the optimisation problem \eqref{eq:error_covariance} under the constraint \eqref{eq:ordinary_kriging_constraint}, and $\xi\in\mathcal{M}_{d_z,d_z}(\mathbb{R})$ the corresponding Lagrange multiplier.

This problem is solved by finding the minimum of $\mathcal{L}$, \textit{i.e.} finding a zero of its derivative. Writing $G$ for $G(\tilde{\mathbf{x}})$, the same properties used in the proof of \Cref{th:simple_kriging} give:
\begin{equation}
	\begin{aligned}
	&\frac{\partial\mathcal{L}}{\partial W}(W,\xi)=0\\
	\iff\quad&2G\:W-2k(\tilde{\mathbf{x}},x)+2\frac{\partial}{\partial W}\left(\operatorname{tr}\left(\xi^\top\left(C^\top W-I_{d_z}\right)\right)\right)=0\\
	\iff\quad&2G\:W-2k(\tilde{\mathbf{x}},x)+2C\xi=0\\
	\iff\quad&G\:W+C\xi=k(\tilde{\mathbf{x}},x)
	\end{aligned}\label{eq:Lagrangian_minimum}
\end{equation}

Together with the constraint $C^\top W=I_{d_z}$, this yields the block system:
\begin{equation}
	\begin{bmatrix}
	G&C\\C^\top&0
	\end{bmatrix}\begin{bmatrix}
	W\\\xi
	\end{bmatrix}=\begin{bmatrix}
	k(\tilde{\mathbf{x}},x)\\I_{d_z}
	\end{bmatrix}
\end{equation}

The left matrix block can be inverted using the upper left Schur inversion formula from \Cref{th:Schur}. Its Schur complement,
\begin{equation}
	S\triangleq0-C^\top G^{-1}C=-C^\top G^{-1}C,
\end{equation}
is invertible since $G\succ0$ and $C$ has full column rank. The convention of \Cref{th:Schur} gives the four blocks of the inverse:
\begin{equation}
	\left\{\begin{aligned}
	\bar{A}&=G^{-1}+G^{-1}C\:S^{-1}C^\top G^{-1}\\
	\bar{B}&=-G^{-1}C\:S^{-1}\\
	\bar{C}&=-S^{-1}C^\top G^{-1}\\
	\bar{D}&=S^{-1}
	\end{aligned}\right.
\end{equation}

Let $M\triangleq-S^{-1}=\left(C^\top G^{-1}C\right)^{-1}$. The weight matrix can be derived from the upper blocks of the inverse:
\begin{equation}
	\begin{aligned}
	W&=\bar{A}\:k(\tilde{\mathbf{x}},x)+\bar{B}\:I_{d_z}\\
	&=\left(G^{-1}+G^{-1}C\:S^{-1}C^\top G^{-1}\right)k(\tilde{\mathbf{x}},x)-G^{-1}C\:S^{-1}\\
	&=G^{-1}k(\tilde{\mathbf{x}},x)+G^{-1}C\:S^{-1}\left(C^\top G^{-1}k(\tilde{\mathbf{x}},x)-I_{d_z}\right)\\
	&=G^{-1}k(\tilde{\mathbf{x}},x)+G^{-1}C\:M\left(I_{d_z}-C^\top G^{-1}k(\tilde{\mathbf{x}},x)\right)
	\end{aligned}
\end{equation}
and $\xi$ from the lower blocks:
\begin{equation}
	\begin{aligned}
	\xi&=\bar{C}\:k(\tilde{\mathbf{x}},x)+\bar{D}\:I_{d_z}\\
	&=-S^{-1}C^\top G^{-1}k(\tilde{\mathbf{x}},x)+S^{-1}\\
	&=S^{-1}\left(I_{d_z}-C^\top G^{-1}k(\tilde{\mathbf{x}},x)\right)\\
	&=-M\left(I_{d_z}-C^\top G^{-1}k(\tilde{\mathbf{x}},x)\right)
	\end{aligned}
\end{equation}

Since $M$ is symmetric, the OK estimator is:
\begin{equation}
	\begin{aligned}
	\hat{z}_{\mathrm{OK}}&=W^\top\tilde{\mathbf{z}}\\
	&=k(x,\tilde{\mathbf{x}})\:G^{-1}\tilde{\mathbf{z}}+\left(I_{d_z}-k(x,\tilde{\mathbf{x}})\:G^{-1}C\right)M\:C^\top G^{-1}\tilde{\mathbf{z}}\\
	&=M\:C^\top G^{-1}\tilde{\mathbf{z}}+k(x,\tilde{\mathbf{x}})\:G^{-1}\left(\tilde{\mathbf{z}}-C\:M\:C^\top G^{-1}\tilde{\mathbf{z}}\right)
	\end{aligned}
\end{equation}
which is exactly \eqref{eq:ordinary_kriging_equations}, the estimated mean being $\hat{m}_{\mathrm{OK}}(\tilde{\mathbf{x}},\tilde{\mathbf{z}})=M\:C^\top G^{-1}\tilde{\mathbf{z}}$.

Since $R_{\mathrm{OK}}$ is a centred moment of the estimation error, it is insensitive to the unknown mean and \eqref{eq:generic_kriging_variance} can be applied. Using \eqref{eq:Lagrangian_minimum} by $W^\top$ and the constraint $W^\top C=I_{d_z}$ gives the identity:
\begin{equation}
	\begin{aligned}
	W^\top G\:W&=W^\top\left(k(\tilde{\mathbf{x}},x)-C\:\xi\right)\\
	&=W^\top k(\tilde{\mathbf{x}},x)-W^\top C\:\xi\\
	&=W^\top k(\tilde{\mathbf{x}},x)-\xi
	\end{aligned}\label{eq:OK_identity}
\end{equation}

Let $U\triangleq I_{d_z}-C^\top G^{-1}k(\tilde{\mathbf{x}},x)$ so that the expressions of $W$ and $\xi$ are simplified to:
\begin{equation}
	W=G^{-1}k(\tilde{\mathbf{x}},x)+G^{-1}C\:M\:U\quad\text{and}\quad\xi=-M\:U
\end{equation}

Combining these results yields:
\begin{equation}
	\begin{aligned}
	R_{\mathrm{OK}}&=W^\top G\:W-W^\top k(\tilde{\mathbf{x}},x)-k(x,\tilde{\mathbf{x}})\:W+k(x,x)+R\\
	&=W^\top k(\tilde{\mathbf{x}},x)-\xi-W^\top k(\tilde{\mathbf{x}},x)-k(x,\tilde{\mathbf{x}})\:W+k(x,x)+R\\
	&=-\xi-k(x,\tilde{\mathbf{x}})\:W+k(x,x)+R\\
	&=M\:U-k(x,\tilde{\mathbf{x}})\:G^{-1}k(\tilde{\mathbf{x}},x)-k(x,\tilde{\mathbf{x}})\:G^{-1}C\:M\:U+k(x,x)+R\\
	&=\left(I_{d_z}-k(x,\tilde{\mathbf{x}})\:G^{-1}C\right)M\:U+k(x,x)+R-k(x,\tilde{\mathbf{x}})\:G^{-1}k(\tilde{\mathbf{x}},x)\\
	&=U^\top M\:U+k(x,x)+R-k(x,\tilde{\mathbf{x}})\:G^{-1}k(\tilde{\mathbf{x}},x)
	\end{aligned}
\end{equation}

Finally, recognising $R_{\hat{m}_{\mathrm{OK}}}(x,\tilde{\mathbf{x}})=U^\top M\:U\succeq0$, since $M\succ0$, gives the desired expression.

\section[Proof of Proposition 3.9]{Proof of \Cref{th:universal_kriging}}\label{sec:proof_UK}

The unbiased assumption now reads:
\begin{equation}
	g(\tilde{\mathbf{x}})^\top W-g(x)^\top=0
\end{equation}
and the Lagrange multiplier method of \Cref{th:ordinary_kriging} can be transposed to universal kriging under the substitutions:
\begin{equation}
	\left\{\begin{aligned}
	C&\longrightarrow g(\tilde{\mathbf{x}})\\
	I_{d_z}&\longrightarrow g(x)^\top\\
	\xi&\in\mathcal{M}_{r_{\mathrm{UK}}d_z,d_z}(\mathbb{R})
	\end{aligned}\right.
\end{equation}

The full column rank of $g(\tilde{\mathbf{x}})$ ensures that $S$ invertible, and that $M$ is SPD. Writing $G$ for $G(\tilde{\mathbf{x}})$, the shorthands of that proof thus become:
\begin{equation}
	\left\{\begin{aligned}
	M&\triangleq-S^{-1}=\left(g(\tilde{\mathbf{x}})^\top G^{-1}g(\tilde{\mathbf{x}})\right)^{-1}\\
	U&\triangleq g(x)^\top-g(\tilde{\mathbf{x}})^\top G^{-1}k(\tilde{\mathbf{x}},x)
	\end{aligned}\right.
\end{equation}
so that the weight matrix and Lagrange multiplier are:
\begin{equation}
	\left\{\begin{aligned}
	W&=G^{-1}k(\tilde{\mathbf{x}},x)+G^{-1}g(\tilde{\mathbf{x}})\:M\:U\\
	\xi&=-M\:U
	\end{aligned}\right.
\end{equation}

Finally, the bilinear identity \eqref{eq:OK_identity} becomes:
\begin{equation}
	W^\top G\:W=W^\top k(\tilde{\mathbf{x}},x)-g(x)\:\xi
\end{equation}

The UK estimator and covariance matrices are derived exactly like for \Cref{th:ordinary_kriging}, using the corrected results.

\section[Proof of Proposition 3.10]{Proof of \Cref{th:kriging_vague_prior}}\label{sec:proof_vague_prior}

\textbf{TODO: à relire}
Let $B\in\mathcal{S}_{r_{\mathrm{UK}}d_z}^{++}(\mathbb{R})$, and give the unknown coefficients the prior $\beta\sim\mathcal{N}(0,B)$, independent of the centred process $f\sim\mathcal{GP}(0,k)$. Randomising the trend this way makes the field $f_B\triangleq f+g\:\beta$ a centred GP of kernel $k_B$ by \Cref{th:GP_marginals}, so that kriging $f_B$ is \Cref{th:simple_kriging} read with $k_B$ in place of $k$. We write $\hat{z}_B$ and $R_B$ for the estimator and covariance it returns, $G$ for $G(\tilde{\mathbf{x}})$ and $g$ for $g(\tilde{\mathbf{x}})$, and reuse the shorthands $M$ and $U$ of \cref{sec:proof_UK}. The Gram matrix of $f_B$ is:
\begin{equation}
	G_B\triangleq k_B(\tilde{\mathbf{x}},\tilde{\mathbf{x}})+I_{\tilde{N}}\otimes\tilde{R}=G+g\:B\:g^\top
\end{equation}

By \Cref{th:Woodbury} and $g^\top G^{-1}g=M^{-1}$, it comes that:
\begin{equation}
	G_B^{-1}=G^{-1}-G^{-1}g\:M_B\:g^\top G^{-1},\qquad M_B\triangleq\left(B^{-1}+M^{-1}\right)^{-1}
\end{equation}

Since $M_B\left(B^{-1}+M^{-1}\right)=I$, we have $M^{-1}M_B=I-B^{-1}M_B$, and thus the two identities:
\begin{equation}
	B-B\:M^{-1}M_B=M_B,\qquad M_B\:M^{-1}B=B-M_B\label{eq:vague_prior_identity}
\end{equation}

Let $A\triangleq k_B(x,\tilde{\mathbf{x}})=k(x,\tilde{\mathbf{x}})+g(x)\:B\:g^\top$. Developing $A\:G_B^{-1}$ and collecting its terms in $B$ through the first identity of \eqref{eq:vague_prior_identity} yields:
\begin{equation}
	A\:G_B^{-1}=k(x,\tilde{\mathbf{x}})\:G^{-1}+U^\top M_B\:g^\top G^{-1}\label{eq:vague_prior_weights}
\end{equation}
whence the estimator:
\begin{equation}
	\hat{z}_B=A\:G_B^{-1}\tilde{\mathbf{z}}=k(x,\tilde{\mathbf{x}})\:G^{-1}\tilde{\mathbf{z}}+U^\top M_B\:g^\top G^{-1}\tilde{\mathbf{z}}
\end{equation}

For the covariance, \eqref{eq:vague_prior_weights} and the second identity of \eqref{eq:vague_prior_identity} give:
\begin{equation}
	\begin{aligned}
	A\:G_B^{-1}A^\top&=k(x,\tilde{\mathbf{x}})\:G^{-1}k(\tilde{\mathbf{x}},x)+U^\top M_B\:g^\top G^{-1}k(\tilde{\mathbf{x}},x)\\
	&\qquad+k(x,\tilde{\mathbf{x}})\:G^{-1}g\:B\:g(x)^\top+U^\top\left(B-M_B\right)g(x)^\top
	\end{aligned}
\end{equation}
and $R_B=k_B(x,x)+R-A\:G_B^{-1}A^\top$ with $k_B(x,x)=k(x,x)+g(x)\:B\:g(x)^\top$. The terms in $B$ cancel, since:
\begin{equation}
	g(x)\:B\:g(x)^\top-k(x,\tilde{\mathbf{x}})\:G^{-1}g\:B\:g(x)^\top-U^\top B\:g(x)^\top=U^\top B\:g(x)^\top-U^\top B\:g(x)^\top=0
\end{equation}
while the two terms in $M_B$ assemble into the sandwich of \eqref{eq:kriged_universal_mean}:
\begin{equation}
	U^\top M_B\:g(x)^\top-U^\top M_B\:g^\top G^{-1}k(\tilde{\mathbf{x}},x)=U^\top M_B\:U
\end{equation}
so that:
\begin{equation}
	R_B=k(x,x)+R-k(x,\tilde{\mathbf{x}})\:G^{-1}k(\tilde{\mathbf{x}},x)+U^\top M_B\:U
\end{equation}

Both expressions are those of \Cref{th:universal_kriging} with $M_B$ substituted for $M$, which proves the stronger statement that kriging under $k_B$ is universal kriging regularised by the prior. Since $M_B\to M$ as $B^{-1}\to0$, the limit returns \eqref{eq:universal_kriging_equations} and \eqref{eq:kriged_universal_mean}.

The coefficients follow the same course. The Gaussian linear model gives $\beta$ the posterior $\mathcal{N}\left(M_B\:g^\top G^{-1}\tilde{\mathbf{z}},M_B\right)$, whose limit is $\mathcal{N}\left(\hat{m}_{\mathrm{UK}}(\tilde{\mathbf{x}},\tilde{\mathbf{z}},g),M\right)$. The estimator of \eqref{eq:kriged_universal_mean} is therefore the posterior mean of the trend, and $M$ its posterior covariance.

\section[Proof of Proposition 3.11]{Proof of \Cref{th:covariance_operator}}\label{sec:proof_Mercer_operator}

As $\mathcal{U}$ is compact and $k$ continuous, $k$ is bounded on $\mathcal{U}\times\mathcal{U}$, and $\mathcal{K}$ is compact.

Let $\phi,\phi'\in\mathcal{L}^2(\mathcal{U})$. Cauchy-Schwarz inequality ensures that there exists a constant $c\in\mathbb{R}_+$ such that:
\begin{equation}
	\iint_{\mathcal{U}\times\mathcal{U}}{\left|\phi(y)^\top k(x,y)^\top\phi'(x)\right|\:dx\:dy}\leq c\:\sqrt{\langle\phi|\phi\rangle}\:\sqrt{\langle\phi'|\phi'\rangle}<+\infty
\end{equation}

Fubini's integration theorem can thus be used to give:
\begin{equation}
	\begin{aligned}
	\langle\mathcal{K}\phi|\phi'\rangle&=\int_{\mathcal{U}}{(\mathcal{K}\phi)(x)^\top\phi'(x)\:dx}\\
	&=\int_{\mathcal{U}}{\left(\int_{\mathcal{U}}{k(x,y)\:\phi(y)\:dy}\right)^\top\phi'(x)\:dx}\\
	&=\int_{\mathcal{U}}{\int_{\mathcal{U}}{\phi(y)^\top k(x,y)^\top\phi'(x)\:dy}\:dx}\\
	&=\int_{\mathcal{U}}{\int_{\mathcal{U}}{\phi(y)^\top k(y,x)\:\phi'(x)\:dx}\:dy}\\
	&=\int_{\mathcal{U}}{\phi(y)^\top\left(\int_{\mathcal{U}}{k(y,x)\:\phi'(x)\:dx}\right)dy}\\
	&=\int_{\mathcal{U}}{\phi(y)^\top(\mathcal{K}\phi')(y)\:dy}\\
	&=\langle\phi|\mathcal{K}\phi'\rangle
	\end{aligned}
\end{equation}

Let $\phi:\mathcal{U}\to\mathbb{R}^{d_z}$ be a continuous function, and let $n\in\mathbb{N}^*$. We define $\left(\mathcal{U}_i^{(n)}\right)_{i\in[\![1,n]\!]}$ a partition of $\mathcal{U}$ such that the diameter of each $\mathcal{U}_i^{(n)},\:i\in[\![1,n]\!]$ tends to $0$ when $n\to+\infty$. Here, the diameter of a subset of the Euclidean space $\mathbb{R}^{d_x}$ is defined as the maximum distance between two points of this subset. For $i\in[\![1,n]\!]$, we define $x_i^{(n)}\in\mathcal{U}_i^{(n)}$, and we let:
\begin{equation}
	a_i^{(n)}\triangleq\left(\int_{\mathcal{U}_i^{(n)}}{dx}\right)\phi\!\left(x_i^{(n)}\right)
\end{equation}

By continuity of $\phi$, the function:
\begin{equation}
	(x,y)\mapsto \phi(x)^\top k(x,y)\:\phi(y)
\end{equation}
is continuous on the compact $\mathcal{U}\times\mathcal{U}$, so its Riemann sums converge:
\begin{equation}
	\sum_{i=1}^n\sum_{j=1}^n\left(a_i^{(n)}\right)^\top k\left(x_i^{(n)},x_j^{(n)}\right)a_j^{(n)}\underset{n\to+\infty}{\longrightarrow}\iint_{\mathcal{U}\times\mathcal{U}}{\phi(x)^\top k(x,y)\:\phi(y)\:dx\:dy}=\langle\phi|\mathcal{K}\phi\rangle
\end{equation}

The SPSDness of $k$ ensures that each Riemann sum is non-negative. By passing to the limit, it follows that:
\begin{equation}
	\langle\phi|\mathcal{K}\phi\rangle\geq0\label{eq:fKf_geq0}
\end{equation}
for every continuous function $\phi$.

Finally, recall that the set of continuous functions of $\mathcal{U}$ is dense in $\mathcal{L}^2(\mathcal{U})$. By continuity of:
\begin{equation}
	\phi\in\mathcal{L}^2(\mathcal{U})\mapsto\langle\phi|\mathcal{K}\phi\rangle
\end{equation}
the inequality \eqref{eq:fKf_geq0} extends to all $\phi\in\mathcal{L}^2(\mathcal{U})$.

\section[Proof of Proposition 3.12]{Proof of \Cref{th:Mercer}}\label{sec:proof_Mercer}

By \Cref{th:covariance_operator}, $\mathcal{K}$ is a compact and self-adjoint operator on $\mathcal{L}^2(\mathcal{U})$. The spectral theorem \cite{bib:spectral-theorem} thus states that there exists an orthonormal basis $(\psi_i)_{i\in\mathbb{N}^*}$ of $\mathcal{L}^2(\mathcal{U})$ consisting of eigenfunctions of $\mathcal{K}$, with each corresponding eigenvalue $(\lambda_i)_{i\in\mathbb{N}^*}$ being real. Up to reordering, the sequence of eigenvalues can be considered to be non-increasing.

Since $k$ is uniformly continuous on the compact $\mathcal{U}\times\mathcal{U}$, $\mathcal{K}\phi$ is continuous for every $\phi\in\mathcal{L}^2(\mathcal{U})$. In particular, for each $i\in\mathbb{N}^*$ such that $\lambda_i\neq0$, the eigenfunction:
\begin{equation}
	\psi_i=\lambda_i^{-1}\mathcal{K}\psi_i
\end{equation}
is continuous. Without loss of generality, it must be noted that the $\psi_i$ associated with a null eigenvalue do not contribute to \eqref{eq:Mercer} and may thus be dismissed.

Since $(\psi_i)_{i\in\mathbb{N}^*}$ is an orthonormal basis of $\mathcal{L}^2(\mathcal{U})$, the family $(\psi_i\:\psi_j^\top)_{i,j\in\mathbb{N}^*}$ is an orthonormal basis of $\mathcal{L}^2(\mathcal{U}\times\mathcal{U})$. Let $\langle\cdot|\cdot\rangle_{\mathcal{L}^2(\mathcal{U}\times\mathcal{U})}$ denote:
\begin{equation}
	\forall\phi,\phi'\in\mathcal{L}^2(\mathcal{U}\times\mathcal{U})\mapsto\langle\phi|\phi'\rangle_{\mathcal{L}^2(\mathcal{U}\times\mathcal{U})}\triangleq\iint_{\mathcal{U}\times\mathcal{U}}{\operatorname{tr}\left(\phi(x,y)^\top\phi'(x,y)\right)\:dx\:dy}
\end{equation}

Let $i,j\in\mathbb{N}^*$. We have:
\begin{equation}
	\begin{aligned}
		\langle k|\psi_i\:\psi_j^\top\rangle_{\mathcal{L}^2(\mathcal{U}\times\mathcal{U})}&=\iint_{\mathcal{U}\times\mathcal{U}}{\operatorname{tr}\left(k(x,y)^\top\psi_i(x)\:\psi_j(y)^\top\right)dx\:dy}\\
		&=\iint_{\mathcal{U}\times\mathcal{U}}{\psi_j(y)^\top k(x,y)^\top\psi_i(x)\:dx\:dy}\\
		&=\int_{\mathcal{U}}{\psi_j(y)^\top\left(\int_{\mathcal{U}}{k(y,x)\:\psi_i(x)\:dx}\right)dy}\\
		&=\int_{\mathcal{U}}{\psi_j(y)^\top\mathcal{K}\psi_i(y)\:dy}\\
		&=\lambda_i\int_{\mathcal{U}}{\psi_j(y)^\top\:\psi_i(y)\:dy}\\
		&=\lambda_i\:\delta_{i,j}
	\end{aligned}
\end{equation}
where:
\begin{equation}
	\delta_{i,j}\triangleq\left\{\begin{aligned}
	1&\quad\text{if }i=j\\0&\quad\text{if }i\neq j
	\end{aligned}\right.
\end{equation}
is the Kronecker symbol.

Because $k\in\mathcal{L}^2(\mathcal{U}\times\mathcal{U})$, it follows from the completeness of $(\psi_i\:\psi_j^\top)_{i,j\in\mathbb{N}^*}$ that:
\begin{equation}
	k_{r_{\mathrm{KL}}}(x,x')\triangleq\sum_{i=1}^{r_{\mathrm{KL}}}\lambda_i\:\psi_i(x)\:\psi_i(x')^\top
\end{equation}
converges to $k$ in the $\mathcal{L}^2(\mathcal{U}\times\mathcal{U})$ sense.

It must be noted that $\mathcal{L}^2$ convergence only identifies $k$ up to a Lebesgue-null subset of $\mathcal{U}\times\mathcal{U}$, and gives no control at any specific point of the domain. In particular, it is not enough to justify substituting $k_{r_{\mathrm{KL}}}$ for $k$ in kriging scenarios. However, since $k$ is continuous on $\mathcal{U}\times\mathcal{U}$, the convergence can be extended to absolute and uniform convergence on the entire domain \cite{bib:Karhunen-Loeve}. The argument applies Dini's theorem \cite{bib:Dini} to the non-decreasing sequence $\left(x\mapsto\operatorname{tr}\:k_{r_{\mathrm{KL}}}(x,x)\right)_{r_{\mathrm{KL}}\in\mathbb{N}^*}$, whose limit is continuous on the compact $\mathcal{U}$, and extends the result off the diagonal through the Cauchy-Schwarz inequality.

Since $\mathcal{K}$ is SPSD, every eigenvalue is non-negative. Assume that there exists a rank $r_0\in\mathbb{N}^*$ such that $\lambda_i=0$ as soon as $i>r_0$. Then the decomposition \eqref{eq:Mercer} would reduce to the finite sum:
\begin{equation}
	k=\sum_{i=1}^{r_0}\lambda_i\:\psi_i\:\psi_i^\top
\end{equation}
whose Gram matrices have rank at most $r_0$ and are thus singular at any $n$ distinct points with $nd_z>r_0$, contradicting the SPDness of $k$ by \Cref{th:kernel_gram_equivalence}. Consequently:
\begin{equation}
	\lambda_1\geq\lambda_2\geq\cdots>0
\end{equation}

\section[Proof of Proposition 3.14]{Proof of \Cref{th:reduced_rank_kriging}}\label{sec:proof_reduced_rank_kriging}

As \eqref{eq:reduced_rank_Gram} is a rank $r_{\mathrm{KL}}$ perturbation of $A\triangleq I_{\tilde{N}}\otimes\tilde{R}$, the Woodbury inversion formula (see \Cref{th:Woodbury}) can be used to reduce the $\tilde{N}d_z\times\tilde{N}d_z$ inversion of $G(\tilde{\mathbf{x}})$ to an $r\times r$ one:
\begin{equation}
	G(\tilde{\mathbf{x}})^{-1}\approx A^{-1}-A^{-1}\Psi(\tilde{\mathbf{x}})\:\tilde{G}(\tilde{\mathbf{x}})^{-1}\Psi(\tilde{\mathbf{x}})^\top A^{-1}
\end{equation}

Computing the centred SK estimator from \Cref{th:simple_kriging} using this approximation and the push-through identity (see \Cref{th:push_through}) yields:
\begin{equation}
	\begin{aligned}
	\hat{z}_{\mathrm{SK}}(x,\tilde{\mathbf{x}},\tilde{\mathbf{z}},0)&=k(x,\tilde{\mathbf{x}})\:G(\tilde{\mathbf{x}})^{-1}\tilde{\mathbf{z}}\\
	&\approx\Psi(x)\:\Lambda\:\Psi(\tilde{\mathbf{x}})^\top\left(I-A^{-1}\Psi(\tilde{\mathbf{x}})\:\tilde{G}(\tilde{\mathbf{x}})^{-1}\Psi(\tilde{\mathbf{x}})^\top\right)A^{-1}\tilde{\mathbf{z}}\\
	&=\Psi(x)\:\Lambda\left(I-\Psi(\tilde{\mathbf{x}})^\top A^{-1}\Psi(\tilde{\mathbf{x}})\:\tilde{G}(\tilde{\mathbf{x}})^{-1}\right)\Psi(\tilde{\mathbf{x}})^\top A^{-1}\tilde{\mathbf{z}}
	\end{aligned}\label{eq:z_s_up_to_tilde_G}
\end{equation}

Recall that, by definition of $\tilde{G}(\tilde{\mathbf{x}})$:
\begin{equation}
	\tilde{G}(\tilde{\mathbf{x}})-\Lambda^{-1}=\Psi(\tilde{\mathbf{x}})^\top A^{-1}\Psi(\tilde{\mathbf{x}})
\end{equation}
which gives:
\begin{equation}
	\begin{aligned}
	I-\Psi(\tilde{\mathbf{x}})^\top A^{-1}\Psi(\tilde{\mathbf{x}})\:\tilde{G}(\tilde{\mathbf{x}})^{-1}&=I-\left(\tilde{G}(\tilde{\mathbf{x}})-\Lambda^{-1}\right)\tilde{G}(\tilde{\mathbf{x}})^{-1}\\
	&=I-I+\Lambda^{-1}\tilde{G}(\tilde{\mathbf{x}})^{-1}\\
	&=\Lambda^{-1}\tilde{G}(\tilde{\mathbf{x}})^{-1}
	\end{aligned}\label{eq:reduced_rank_lemma}
\end{equation}

Plugging this result into \eqref{eq:z_s_up_to_tilde_G} and using $A^{-1}=I_{\tilde{N}}\otimes\tilde{R}^{-1}$ finally gives:
\begin{equation}
	\hat{z}_{\mathrm{SK}}(x,\tilde{\mathbf{x}},\tilde{\mathbf{z}},0)\approx\Psi(x)\:\tilde{G}(\tilde{\mathbf{x}})^{-1}\Psi(\tilde{\mathbf{x}})^\top A^{-1}\tilde{\mathbf{z}}
\end{equation}

For $R_{\mathrm{SK}}$, reusing:
\begin{equation}
	k(x,\tilde{\mathbf{x}})\:G(\tilde{\mathbf{x}})^{-1}\approx\Psi(x)\:\tilde{G}(\tilde{\mathbf{x}})^{-1}\Psi(\tilde{\mathbf{x}})^\top A^{-1}
\end{equation}
together with \eqref{eq:reduced_rank_lemma} yields:
\begin{equation}
	\begin{aligned}
	R_{\mathrm{SK}}&=k(x,x)+R-k(x,\tilde{\mathbf{x}})\:G(\tilde{\mathbf{x}})^{-1}k(\tilde{\mathbf{x}},x)\\
	&\approx\Psi(x)\:\Lambda\:\Psi(x)^\top+R-\Psi(x)\:\tilde{G}(\tilde{\mathbf{x}})^{-1}\Psi(\tilde{\mathbf{x}})^\top A^{-1}\times\Psi(\tilde{\mathbf{x}})\:\Lambda\:\Psi(x)^\top\\
	&=\Psi(x)\:\Lambda\:\Psi(x)^\top+R-\Psi(x)\:\tilde{G}(\tilde{\mathbf{x}})^{-1}\left(\tilde{G}(\tilde{\mathbf{x}})-\Lambda^{-1}\right)\Lambda\:\Psi(x)^\top\\
	&=\Psi(x)\:\Lambda\:\Psi(x)^\top+R-\Psi(x)\left(\Lambda-\tilde{G}(\tilde{\mathbf{x}})^{-1}\right)\Psi(x)^\top\\
	&=R+\Psi(x)\:\tilde{G}(\tilde{\mathbf{x}})^{-1}\Psi(x)^\top
	\end{aligned}
\end{equation}

Finally, when $\tilde{R}=\sigma^2I_{d_z}$ with $\sigma>0$, substituting $A^{-1}=\sigma^{-2}I_{\tilde{N}d_z}$ yields the announced simplifications. For such $\sigma>0$, the matrix:
\begin{equation}
	\Psi(\tilde{\mathbf{x}})^\top\Psi(\tilde{\mathbf{x}})+\sigma^2\Lambda^{-1}\succ0
\end{equation}
is SPD, as the sum of the SPSD matrix $\Psi(\tilde{\mathbf{x}})^\top\Psi(\tilde{\mathbf{x}})$ and the SPD matrix $\sigma^2\Lambda^{-1}$. The map:
\begin{equation}
	\sigma\in\mathbb{R}_+^*\mapsto\left(\Psi(\tilde{\mathbf{x}})^\top\Psi(\tilde{\mathbf{x}})+\sigma^2\Lambda^{-1}\right)^{-1}
\end{equation}
is thus well-defined and continuous, and extends to $\sigma=0$ when $\Psi(\tilde{\mathbf{x}})^\top\Psi(\tilde{\mathbf{x}})\succ0$, \textit{i.e.} when $\Psi(\tilde{\mathbf{x}})$ has full column rank. As $\hat{z}_{\mathrm{SK}}$ and $R_{\mathrm{SK}}$ depend on $\sigma$ only through this inverse and an explicit factor $\sigma^2$, both expressions remain valid at $\sigma=0$ under this condition.

\section[Proof of Proposition 3.16]{Proof of \Cref{th:native_RKHS}}\label{sec:proof_native_RKHS}

\textbf{TODO: à relire}

\textbf{The form is well defined.} Let $\mathcal{H}_0$ denote the span of $\{k(x,\cdot),\:x\in\mathcal{U}\}$, and extend $\langle\cdot|\cdot\rangle_{\mathcal{H}_k}$ to it by bilinearity. For $\psi=\sum_{i=1}^na_i\:k(x_i,\cdot)$ and $\varphi=\sum_{j=1}^pb_j\:k(x'_j,\cdot)$ in $\mathcal{H}_0$, it comes that:
\begin{equation}
	\langle\psi|\varphi\rangle_{\mathcal{H}_k}=\sum_{i=1}^n\sum_{j=1}^pa_i\:b_j\:k(x_i,x'_j)=\sum_{j=1}^pb_j\:\psi(x'_j)=\sum_{i=1}^na_i\:\varphi(x_i)\label{eq:RKHS_bilinear}
\end{equation}
The last two expressions involve one of the two functions through its values alone, so neither representation matters and the form is well defined. It is symmetric and bilinear by construction, and positive since $k$ is SPSD. Taking $\varphi=k(x,\cdot)$ gives the reproducing property on $\mathcal{H}_0$:
\begin{equation}
	\forall\psi\in\mathcal{H}_0,\:\forall x\in\mathcal{U},\:\langle\psi|k(x,\cdot)\rangle_{\mathcal{H}_k}=\psi(x)\label{eq:reproducing_generators}
\end{equation}

\textbf{It is definite.} The Cauchy-Schwarz inequality holds for any positive semi-definite bilinear form, so \eqref{eq:reproducing_generators} bounds every evaluation:
\begin{equation}
	\forall\psi\in\mathcal{H}_0,\:\forall x\in\mathcal{U},\:|\psi(x)|^2\leq\langle\psi|\psi\rangle_{\mathcal{H}_k}\:k(x,x)\label{eq:RKHS_evaluation_bound}
\end{equation}
A function of null norm thus vanishes everywhere, and $\langle\cdot|\cdot\rangle_{\mathcal{H}_k}$ is an inner product on $\mathcal{H}_0$.

\textbf{The completion is a space of functions.} \Cref{eq:RKHS_evaluation_bound} makes every Cauchy sequence of $\mathcal{H}_0$ pointwise Cauchy, hence pointwise convergent, which sends each class of the abstract completion to a function on $\mathcal{U}$. That map is injective. Indeed, let $(\psi_n)_n$ be a Cauchy sequence converging pointwise to $0$, let $\varepsilon>0$, and let $N$ be such that $\|\psi_n-\psi_m\|_{\mathcal{H}_k}\leq\varepsilon$ for $n,m\geq N$. Fixing $n\geq N$ and writing $\psi_n=\sum_{i=1}^{n'}a_i\:k(x_i,\cdot)$, \eqref{eq:RKHS_bilinear} gives:
\begin{equation}
	\|\psi_n\|_{\mathcal{H}_k}^2=\langle\psi_n|\psi_n-\psi_m\rangle_{\mathcal{H}_k}+\sum_{i=1}^{n'}a_i\:\psi_m(x_i)
\end{equation}
whose first term is at most $\varepsilon\sup_n\|\psi_n\|_{\mathcal{H}_k}$, finite since a Cauchy sequence is bounded, and whose second is a finite sum tending to $0$ with $m$. Hence $\|\psi_n\|_{\mathcal{H}_k}\to0$, and only the zero class is sent to the zero function. Passing \eqref{eq:reproducing_generators} to the limit, $\mathcal{H}_k$ is an RKHS of reproducing kernel $k$.

\textbf{Uniqueness.} Let $\mathcal{H}$ be an RKHS of reproducing kernel $k$. It contains every $k(x,\cdot)$, hence $\mathcal{H}_0$, and its inner product agrees with $\langle\cdot|\cdot\rangle_{\mathcal{H}_k}$ on the generators, hence on $\mathcal{H}_0$ by bilinearity. A function of $\mathcal{H}$ orthogonal to $\mathcal{H}_0$ verifies $\psi(x)=\langle\psi|k(x,\cdot)\rangle_{\mathcal{H}}=0$ for every $x\in\mathcal{U}$, so $\mathcal{H}_0$ is dense in $\mathcal{H}$, which is therefore its completion $\mathcal{H}_k$.
